\documentclass[../preamble]{subfiles}
\begin{document}

\section{Natural Isomorphisms}
\label{sec:natural-isomorphism}
\concept{natural-isomorphism}{Natural Isomorphism}
\depends{natural-transformation}
\implements{haskell}{Cat.NaturalIsomorphism}
\implements{agda}{Cat.NaturalIsomorphism}
\implements{lisp}{natural-isomorphism}

Let $\Category{C}$ and $\Category{D}$ be categories,
and let $\Functor{F}, \Functor{G} \colon \Category{C} \lto \Category{D}$ be functors.
A \newterm{natural isomorphism}
$\newmath{\NatTrans{\alpha} \colon \Functor{F} \xRightarrow{\sim} \Functor{G}}$
is a natural transformation $\NatTrans{\alpha} \colon \Functor{F} \Rightarrow \Functor{G}$
such that for every object $X \in \Ob(\Category{C})$,
the component $\NatTrans{\alpha}_X \colon \Functor{F}(X) \lto \Functor{G}(X)$
is an isomorphism in $\Category{D}$
(\nlabref{natural+isomorphism}).

\begin{definition}[Natural isomorphism, componentwise]
  \label{def:nat-iso-componentwise}
  Equivalently,
  a natural isomorphism $\Functor{F} \cong \Functor{G}$
  consists of a pair of natural transformations
  $\newmath{\NatTrans{\alpha} \colon \Functor{F} \Rightarrow \Functor{G}}$
  and
  $\newmath{\NatTrans{\alpha}^{-1} \colon \Functor{G} \Rightarrow \Functor{F}}$
  such that their vertical composites
  (Definition~\ref{def:vertical-composition})
  satisfy
  \begin{equation}
    \label{eq:nat-iso-left}
    \NatTrans{\alpha}^{-1} \circ \NatTrans{\alpha} = \id_{\Functor{F}},
  \end{equation}
  \begin{equation}
    \label{eq:nat-iso-right}
    \NatTrans{\alpha} \circ \NatTrans{\alpha}^{-1} = \id_{\Functor{G}}.
  \end{equation}
  At each component,
  these equations assert that $\NatTrans{\alpha}_X$ and $\NatTrans{\alpha}^{-1}_X$
  are mutually inverse morphisms in $\Category{D}$.
\end{definition}

The commutativity of the following diagram expresses
that $\NatTrans{\alpha}$ and $\NatTrans{\alpha}^{-1}$
are mutually inverse at each component:
\begin{equation}
  \label{eq:nat-iso-roundtrip}
  \begin{tikzcd}
    \Functor{F}(X) \arrow[r, "\NatTrans{\alpha}_X", shift left]
      & \Functor{G}(X) \arrow[l, "\NatTrans{\alpha}^{-1}_X", shift left]
  \end{tikzcd}
\end{equation}

\begin{definition}[Identity natural isomorphism]
  \label{def:id-nat-iso}
  The \newterm{identity natural isomorphism}
  $\newmath{\id_{\Functor{F}} \colon \Functor{F} \xRightarrow{\sim} \Functor{F}}$
  has forward and backward components both equal to the identity:
  $(\id_{\Functor{F}})_X = \id_{\Functor{F}(X)}$
  and $(\id_{\Functor{F}})^{-1}_X = \id_{\Functor{F}(X)}$.
  Equations~\eqref{eq:nat-iso-left} and~\eqref{eq:nat-iso-right}
  hold by the unit law.
\end{definition}

\begin{definition}[Composition of natural isomorphisms]
  \label{def:compose-nat-iso}
  Given natural isomorphisms
  $\NatTrans{\alpha} \colon \Functor{F} \xRightarrow{\sim} \Functor{G}$
  and
  $\NatTrans{\beta} \colon \Functor{G} \xRightarrow{\sim} \Functor{H}$,
  their \newterm{composite natural isomorphism}
  $\newmath{\NatTrans{\beta} \circ \NatTrans{\alpha}
    \colon \Functor{F} \xRightarrow{\sim} \Functor{H}}$
  has forward components
  $(\NatTrans{\beta} \circ \NatTrans{\alpha})_X
    = \NatTrans{\beta}_X \circ \NatTrans{\alpha}_X$
  and backward components
  $(\NatTrans{\beta} \circ \NatTrans{\alpha})^{-1}_X
    = \NatTrans{\alpha}^{-1}_X \circ \NatTrans{\beta}^{-1}_X$.
\end{definition}

\begin{definition}[Inverse of a natural isomorphism]
  \label{def:invert-nat-iso}
  Given a natural isomorphism
  $\NatTrans{\alpha} \colon \Functor{F} \xRightarrow{\sim} \Functor{G}$,
  its \newterm{inverse natural isomorphism}
  $\newmath{\NatTrans{\alpha}^{-1}
    \colon \Functor{G} \xRightarrow{\sim} \Functor{F}}$
  is obtained by swapping forward and backward components.
\end{definition}

\begin{remark}
  Natural isomorphisms are the isomorphisms in the functor category
  $[\Category{C}, \Category{D}]$.
  Two functors are \newterm{naturally isomorphic},
  written $\Functor{F} \cong \Functor{G}$,
  when there exists a natural isomorphism between them.
  This is the correct notion of ``sameness'' for functors.
\end{remark}

\end{document}
